\section{Related Work}
\label{sec:related}

\para{LLMs for human-comparable event forecasting.} The most relevant line of work studies binary event forecasting on real questions with human crowd baselines. Halawi \etal develop a retrieval-augmented language-model forecasting pipeline and show that it approaches, but does not generally exceed, human crowd performance on a large curated benchmark from multiple forecasting platforms \cite{halawi2024approaching}. Their strongest gains appear when the crowd is uncertain or when forecasts are made earlier in a question's lifecycle. Schoenegger and Park report a more negative result: GPT-4 underperforms a live Metaculus crowd and does not significantly beat a 0.5 baseline in a tournament-style evaluation \cite{schoenegger2023llm}. Paleka \etal focus on logical consistency checks for language-model forecasters rather than direct human comparison, but their framework is relevant because it offers a way to diagnose forecast quality before outcomes resolve \cite{paleka2024consistency}. Our work sits closest to this literature. Unlike large retrieval systems, we use a controlled local-context setup, and unlike live-tournament studies, we vary information regime within one dataset.

\para{LLMs for numeric time-series forecasting.} A separate literature treats forecasting as sequence extrapolation over dense numeric observations. Gruver \etal show that general language models can perform zero-shot time-series forecasting by serializing numeric sequences into text \cite{gruver2023llmtime}. Jin \etal propose Time-LLM, which reprograms frozen LLMs for time-series forecasting through patch embeddings and prompt-style conditioning \cite{jin2024timellm}. Liu \etal extend this direction with large time-series foundation models trained directly for temporal tasks \cite{liu2024timer}, and Guo \etal argue that temporal distillation can further improve LLM-based forecasting \cite{guo2026tllm}. These papers are important evidence that LLM-style architectures can work well on high-data forecasting problems, but they do not answer whether such systems outperform human crowd judgments on event questions.

\para{Skeptical ablations and task-definition concerns.} Later analyses caution against attributing time-series gains too broadly to language modeling itself. Tan \etal reproduce several influential LLM time-series systems and find that simpler ablations often match or exceed them, suggesting that gains may come from patching or attention structure rather than language transfer \cite{tan2024timeseries}. This critique matters for our question because the user hypothesis mixes two forecasting settings that behave differently in the literature. We therefore stay within binary question forecasting and vary information regime directly instead of combining event forecasting with dense numeric benchmarks. That choice makes our comparison narrower, but it gives a cleaner answer to when an LLM can match or beat a human crowd on the same underlying task.

